% Esquisse d'un programme (cote 119, copy 1): archive cote 119, site page 22 (PDF page 23), typescript page 20.
% typescript. Draft compared: EsquisseFr.pdf (text layer):743-795.
% AI-transcribed by Claude Opus 5.5 (claude-opus-5-5, adaptive thinking, effort high), Message Batches API, batch esq1 (msgbatch_01LNsogFyF3y6KSp9UbtJtBk), 2026-11-12; not yet checked by eye.
\scanpage{119/22}
\mspage{20}

d'une carte combinatoire garde un sens sur un corps de base, et
même sur un anneau de base arbitraire. Elle garde également un
sens pour des polyèdres réguliers de dimension quelconque, en
remplaçant le groupe cartographique $\underline{C}_2$ par une variante n-dimen-%
sionnelle $\underline{C}_n$ convenable. Le cas n = 1, i.e. la théorie des
polygônes réguliers en caractéristique quelconque, fait l'objet
d'un cours DEA en 1977/78, et fait apparaître déjà quelques phé-%
nomènes nouveaux, comme aussi l'utilité de travailler non pas
dans un espace ambiant affine (ici le plan affine), mais dans
un espace \uline{projectif}. Ceci est dû notamment au fait que dans cer-%
taines caractéristiques (et notamment en caractéristique 2) le
centre d'un polyèdre régulier est rejeté à l'infini. D'autre
part, le contexte projectif, contrairement au contexte affine,
permet de développer avec aisance un formalisme de dualité pour
les polyèdres réguliers, correspondant au formalisme de dualité
des cartes combinatoires ou topologiques (où le rôle des sommets
et des faces, dans le cas n = 2 disons, se trouve interchangé).
Il se trouve que pour tout polyèdre régulier projectif, on peut
définir un hyperplan canonique associé, qui joue le rôle d'un
hyperplan à l'infini canonique, et permet de considérer le polyè-%
dre donné comme un polyèdre régulier affine.

L'extension de la théorie des polyèdres réguliers (et
plus généralement, de toutes sortes de configurations géométrico-%
combinatoires, y compris les systèmes de racines...) du corps de
base $\mathbb{R}$ ou $\mathbb{C}$ vers un anneau de base général, me semble d'une
portée comparable, dans cette partie de la géométrie, à l'exten-%
sion analogue qui a eu lieu depuis le début du siècle en géométrie
algébrique, ou depuis une vingtaine d'années en topologie ( *),
avec l'introduction du langage des schémas et celui des topos.

\rule{4cm}{0.4pt}

(*) En écrivant cela, je suis conscient que rares sont les topo-%
logues, encore aujourd'hui, qui se rendent compte de cet élargisse-%
ment conceptuel et technique de la topologie, et des ressources
qu'elle offre.

\margin{22}
